\section{Stock Price Fragility}
\label{sec:stock_price_fragility}

While the fire-sale literature explains how liquidity-driven trading can generate non-fundamental price pressure, it does not by itself identify which securities are most exposed to that pressure. \textcite{greenwood2011stock} address this question by introducing the concept of stock price fragility. They define an asset as fragile when it is susceptible to non-fundamental shifts in demand. Fragility can arise because ownership is concentrated or because the asset's owners face volatile or correlated liquidity shocks that cause them to trade simultaneously. It is therefore a property of the ownership base as a whole rather than of any individual investor. Empirically, the measure combines the composition of ownership with the estimated variance--covariance structure of owners' liquidity-driven trades. \textcite{greenwood2011stock} show that greater fragility strongly predicts higher subsequent total and idiosyncratic return volatility.

This formulation is related to crowding, but it contains more information than ownership concentration alone because it accounts for the volatility and correlation of owners' expected liquidity trades. More recent evidence moves toward the specific downside outcome examined in this thesis. Using hedge-fund position data, \textcite{brown2022crowded} construct security-level measures of crowdedness and show that stocks with greater exposure to crowdedness experience relatively larger drawdowns during periods of market distress. Their findings indicate that ownership-based vulnerability can become particularly visible in adverse market states and provide direct empirical motivation for using future maximum drawdown as the primary downside-risk target.